\documentclass[11pt,a4paper]{article}
\usepackage[T1]{fontenc}
\usepackage[utf8]{inputenc}
\usepackage[polish]{babel}
\usepackage{lmodern,microtype}
\usepackage[a4paper,margin=2.3cm]{geometry}
\usepackage{booktabs,longtable,xcolor,hyperref,listings}
\hypersetup{colorlinks=true,linkcolor=blue,urlcolor=blue}
\lstset{basicstyle=\ttfamily\small,breaklines=true,frame=single}
\title{Etapy 2 i 3 Smog AI\\Model lokalny, aplikacja lokalna i DigitalOcean}
\author{Edward Szczypka}
\date{\today}
\begin{document}
\maketitle

\section{Etap 2}
Etap modelowy obejmuje: snapshot, wybór ograniczonego zbioru, trening quick/full,
monitoring, bramę jakości oraz atomową aktywację. Import historyczny może działać
równolegle, ponieważ model czyta niezmienny dataset.

\begin{lstlisting}[language=bash]
.\scripts\Run-Stage2-PM25-Pilot.ps1 `
  -ProjectRoot $ProjectRoot `
  -RuntimeRoot $RuntimeRoot `
  -Parameters "PM2.5" `
  -Snapshot auto
\end{lstlisting}

\section{Etap 3A --- lokalne serving}
FastAPI i Streamlit nie trenują modeli. Czytają ostatnie poprawne artefakty.
Lokalna bramka obejmuje health, zapytanie tekstowe, profil godzinowy, mapę,
dokumentację i katalog modeli.

\section{Etap 3B --- DigitalOcean App Platform}
Docelowy przepływ:
\[
 \text{Windows ML}\to\text{DigitalOcean Spaces}\to
 \text{FastAPI}\to\text{Streamlit}.
\]
App Platform nie posiada bazy danych ani joba ML. FastAPI otrzymuje klucz Spaces
z uprawnieniem odczytu, a dashboard komunikuje się z FastAPI przez prywatny URL.

\section{Automatyczne wdrożenie}
Workflow GitHub wykonuje testy, waliduje App Spec, wdraża dwa serwisy i wykonuje
smoke test. Push nie może ominąć bramy testowej.

\section{Preflight}
\begin{lstlisting}[language=bash]
.\scripts\Test-Stage2Stage3Readiness.ps1 `
  -ProjectRoot $ProjectRoot `
  -RuntimeRoot $RuntimeRoot `
  -VerifySnapshotChecksum `
  -StrictArtifacts
\end{lstlisting}

\end{document}
